% ------------------------------ define equals
% also \coloneqq and \eqqcolon from mathtools imported below
\newcommand{\defeq}{\vcentcolon=}
\newcommand{\eqdef}{=\vcentcolon}

% ------------------------------ Expectation, Var and Covar symbol
% when declaring math operators the * denotes the use of limits above and/or below
\DeclareMathOperator*{\E}{\mathbb{E}}
\DeclareMathOperator*{\Var}{\operatorname{Var}}
\DeclareMathOperator*{\Cov}{\operatorname{Cov}}


% ----------------------------- set complement
\newcommand{\stcomp}[1]{{#1}^{\mathsf{c}}}
% ----------------------------- argmax and argmin
\DeclareMathOperator*{\argmax}{arg\,max\,}
\DeclareMathOperator*{\argmin}{arg\,min\,}

% ----------------------------- transpose and other matrix notation
% https://tex.stackexchange.com/a/217623/151207
% unicode-math not supported in pdftex
% \usepackage{mathtools} % in preamble
%\usepackage{unicode-math}
\newcommand*{\matr}[1]{\mathbf{#1}}
\newcommand*{\tran}{^{\mkern-1.5mu\mathsf{T}}}
\newcommand*{\conj}[1]{\overline{#1}}
\newcommand*{\hermconj}{^{\mathsf{H}}}
% trace of a matrix
\DeclareMathOperator{\Tr}{Tr}
% cardinality of a matrix and vector
\newcommand{\matdim}[2]{\mathsf{\textit{#1}\times\textit{#2}}}
\newcommand{\vecdim}[1]{\mathsf{\textit{#1}}}

% Define a command for a placeholder matrix with dimensions m x n
\newcommand{\matblock}[3]{
	\begin{bmatrix}
		#1_{\mathsf{\textit{1,1}}}  & \cdots & #1_{\mathsf{\textit{1,#3}}}  \\
		\vdots                      & \ddots & \vdots                       \\
		#1_{\mathsf{\textit{#2,1}}} & \cdots & #1_{\mathsf{\textit{#2,#3}}}
	\end{bmatrix}
}

% Define a command for a placeholder vector with dimension m
\newcommand{\vecblock}[2]{
	\begin{bmatrix}
		#1_{\mathsf{\textit{1}}} \\
		\vdots                   \\
		#1_{\mathsf{\textit{#2}}}
	\end{bmatrix}
}
